\section{Asking the Unchanging God}

A sharper objection now appears. If God knows all things and his providence is not revised by new information, what can petition possibly do? Intercession can seem to require an unstable God: first unwilling, then persuaded; first unaware, then informed.

Christian prayer requires no such change. Aquinas distinguishes the eternal divine disposition from the created means through which its effects occur. Providence ordains not only that a good will be given, but that it will be given through particular causes. Rain nourishes through soil; medicine heals through bodily powers. Such agents are \emph{secondary causes}: real causes whose power and opportunity are themselves received. Prayer can be one of them. We ask to receive, through asking, what God eternally wills to give through that prayer.\endnote{ST II--II, q.~83, a.~2, response and replies 1--3: providence orders effects together with their created causes.}

This does not reduce prayer to an inward exercise. The first thing laid open is the person praying: mind raised, desire named, dependence confessed, will made available. The \emph{Catechism} calls transformation of the praying heart the ``first response'' to petition, not the only response. It then speaks of Christ's prayer as efficacious and Christian prayer as cooperation with providence.\endnote{CCC 2735--2741, especially 2738--2741: inward transformation and efficacious petition are held together.}

Intercession is therefore neither information supplied to omniscience nor a pious fiction. It is charity taking the form of asking. God can eternally will that one person's healing, courage, conversion, protection, or material help be given through another's prayer, as through a physician, friend, judge, or stranger. We often cannot identify the effect or measure our contribution. Hiddenness limits knowledge; it does not make the cause unreal. ``Pray for one another,'' James commands, because ``the prayer of a righteous man has great power in its effects.''\endnote{Jas 5:16; 1 Tim 2:1--4; CCC 2634--2636. No particular temporal outcome is promised.}

Then may one ask for temporal goods? Certainly. Christ teaches, ``Give us this day our daily bread.'' But temporal goods must be desired under God, not as religion's hidden final object. Aquinas permits asking for what supports bodily life and virtue; such things are sought second in importance, while union with God is sought first.\endnote{Matt 6:9--13; ST II--II, q.~83, aa.~5--6. Created goods remain goods, not the price of loving God.}

Prayer therefore carries no guarantee of wealth, relief, marriage, a particular vocation, repaired family relations, emotional serenity, or a decipherable sign. It cannot replace apology, treatment, deliberation, court action, study, employment, or the patient work by which many temporal goods ordinarily come. To promise these outcomes in exchange for a devotion is to turn the Giver into a mechanism and the devotion into leverage. The scandal is not merely that the promise may fail. It teaches the heart to want God as the means to something else.

The final object of prayer is the opposite. The Psalmist says, ``One thing have I asked of the Lord, this will I seek: that I may dwell in the house of the Lord all the days of my life.''\endnote{Ps 27:4 (Ps 26:4 Vulgate); ST II--II, q.~83, a.~1 ad 2. ``Object'' means the good principally sought.} Petition, intercession, thanksgiving, sorrow, and praise become modes of one desire: not to drag God into my plan, but to live every created good within his.
